\begin{frame}{The Vision for Autonomy in Defense}
\begin{columns}
    \begin{column}{0.6\textwidth}
        \begin{itemize}
            \item 
                \textbf{DoD initiatives for autonomous systems.}
            \begin{itemize}
                \item 
                    \textbf{Replicator Initiative:}
                    Calls for the rapid integration of commercial-off-the-shelf (COTS)
                    autonomous systems for scalable deployment. \cite{robertson2023}
                \note[item]{Deputy Sec.\ Hicks announced replicator in 2023}
                \note[item]{Replicator: Reps COTS motivation}
                \note[item]{    Which is both a production economy plus}
                \item 
                    \textbf{DARPA's OFFSET:}
                    Demonstrated feasibility of swarm-enabled ground and aerial 
                    assets in contested environments. \cite{zotero-2835}
                \note[item]{OFFSET: Reps combined arms motivation}
            \end{itemize}
        \end{itemize}
    \end{column}
    \begin{column}{0.4\textwidth}
        \begin{figure}[!h]
            \centering
            \includegraphics[width=0.95\textwidth]{replicator_collage.png}
            \caption{
                \textbf{Replicator Initiative} intended for multiple platforms. 
                \cite{robertson2023}
            }
            \label{fig:replicator_collage}
        \end{figure}
    \end{column}
\end{columns}
\end{frame}


\begin{frame}{The Training Bottleneck}
    \begin{itemize}
        \item \textbf{Sample complexity scales with team size.}
            Learning in multi-agent systems requires exponentially more
            environment interactions as team size grows~\cite{jin2025,canese2021}.

        \item \textbf{Benchmark $\neq$ operational.}
            Controlled evaluation environments mask the cost gap between
            research settings and real-world deployment pipelines
            \cite{gronauer2022,krouka2022}.

        \item \textbf{Current training pipelines do not generalize.}
            Policies trained for one team size or composition
            typically must be retrained from scratch for any
            change in team structure.

        \item \textbf{Bottom line:} Overcoming these barriers requires
            rethinking \emph{how} we represent the learning problem —
            not merely which algorithm we apply.
    \end{itemize}

    \note{This is a transitional slide — keep it brisk (about 1 minute).
    The key message is that the computational bottleneck is \emph{real and structural},
    not just a hardware problem that more GPUs will solve.
    The audience should leave this slide with a clear sense of urgency:
    training cost is the primary reason MARL has not translated from
    benchmark wins to operational fielding.
    The final bullet plants the seed for the representational thesis
    you will state explicitly on slide 1.5.}
\end{frame}

\begin{frame}{Heterogeneity as the Core Challenge}
    \begin{itemize}
        \item Real swarms may be \textbf{non-homogeneous}.
            Operational teams mix UAVs, UGVs, sensors, and roles
            \cite{rizk2019,hoang2023}.

        \item Heterogeneity is a \textbf{requirement} for effectiveness
            and a \textbf{complicating factor} for training
            \cite{calvo2018,zhong2024}.
    \end{itemize}
\begin{columns}
    \begin{column}{0.55\textwidth}

        % \vspace{0.5em}
        \begin{block}{Two Kinds of Heterogeneity}
            \renewcommand{\arraystretch}{1.3}
            \begin{tabular}{@{} p{2.1cm} >{\raggedright\arraybackslash}p{4.9cm} @{}}
                \textbf{Behavioral} &
                    Structurally identical agents developing different roles \\
                \textbf{Intrinsic} &
                    Agents differ in sensors, effectors, or observation spaces \\
            \end{tabular}
        \end{block}

        \vspace{0.3em}
        \small This distinction is not merely taxonomic —
        it \textbf{determines which methods are applicable}.
    \end{column}
    \begin{column}{0.45\textwidth}
        \begin{figure}
            \centering
            % \includegraphics[width=0.95\linewidth]{offset_collage.jpg}
            \includegraphics[width=0.85\textwidth]{hae_cycle.tex}
            \caption{Mixed-platform teams combine agents with
                     different sensors and roles.}
        \end{figure}
    \end{column}
\end{columns}

    \note{Spend about 90 seconds here.
    The behavioral vs.\ intrinsic taxonomy is foundational vocabulary
    for everything that follows — the committee needs to have it before
    you get to the contribution slides.\\
    Emphasize 
    \emph{Behavioral}: as emergent, learned differences
    \emph{Intrinsic}: as structural, based on some attribute(s) of the agents}
    \note[item]{
        Replicator and OFFSET call to the desire for more than a fleet of identical quadcopters.
    }
    \note[item]{
        Call out the asymmetry: the capability advantage of mixed platforms 
        comes at a real training cost.
        But in hardware: a potential for savings,
        like my kitchen utensils when I move out, cheap and mismatched
    }
\end{frame}


\begin{frame}{Why Standard MARL Breaks Down}
    \begin{itemize}
        \item Homogeneous MARL assumes a shared domain and image for the policy function.
        \item Intrinsic heterogeneity implies a difference in domain or image.
        \item Making parameter sharing difficult or impossible.
    \end{itemize}
    \begin{center}
    \renewcommand{\arraystretch}{1.4}
    \footnotesize
    \begin{tabular}{@{} l l l @{}}
        \toprule
        \textbf{Approach} & \textbf{Key Mechanism} & \textbf{Limitation} \\
        \midrule
        Per-type policies~\cite{papoudakis2021} & Separate network per type
            & $|{I}|\times$ storage; no shared learning \\
        One-hot agent IDs~\cite{terry2020} & Append type label to obs.
            & Arbitrary; fails on zero-shot novel agents \\
        GNN / Architectural~\cite{liu2020b} & Graph aggregation of input
            & Added complexity; inductive bias mismatch \\
        \textbf{Implicit Indication} & Homogenized obs.\ space
            & \textbf{Single network, no behavioral-HARL} \\
        \bottomrule
    \end{tabular}
    \end{center}

    \note[item]{MARL typically requires a consistent input (observation) space, and common output (action) space}
    \note[item]{as a function, the policy has one domain and image, and agents are expected to reflect it}

    \note{
        We found approaches to  behavioral-HARL to be more common.\\
        This table highlights the main approaches that we found:
    }
    \note[item]{Most straight-forward: give them all policies}
    \note[item]{Tried to work around this in C1}
    \note[item]{One-hot (Explicit Encoding): provides a flag to dif agents, great for emergent HARL}
    \note[item]{GNN: usually as a means of communication between agents, but can be used in HARL\\
        We found it vulnerable to inductive bias.}
    \note[item]{Implicit Indication proposed in this work: big limit on emergent/behavioral HARL}
    
    \note[item]{In C2 (Diss:Chapter 3) (Def:Section 4) we do II}
    \note[item]{In C3 (Diss:Chapter 4) (Def:Section 5) we eval against a GNN}
    \note{\bfseries TIME:}
\end{frame}


\begin{frame}[shrink=18]{Research Questions}
\vspace{1em}
\begin{center}
\renewcommand{\arraystretch}{1.2}
\small
\begin{tabular}{@{} l >{\raggedright\arraybackslash}p{0.9\linewidth} @{}}
    \toprule
    \multicolumn{2}{@{}l}{\textbf{C1 — Policy Upsampling}} \\
    \midrule
    \textbf{RQ1.1} & Does smaller-team pretraining + policy duplication
                      improve efficiency without sacrificing performance? \\
    \textbf{RQ1.2} & How does effectiveness vary across environments with
                      different heterogeneity forms? \\
    \midrule
    \multicolumn{2}{@{}l}{\textbf{C2 — Implicit Indication}} \\
    \midrule
    \textbf{RQ2.1} & How do invariant observation spaces affect learning
                      efficiency? \\
    \textbf{RQ2.2} & Do they stabilize performance under team-size changes,
                      sensor dropout, and novel compositions? \\
    \textbf{RQ2.3} & What are the computational costs relative to per-type policies? \\
    \midrule
    \multicolumn{2}{@{}l}{\textbf{C3 — Paradigm Comparison}} \\
    \midrule
    \textbf{RQ3.1} & To what extent do graph-based policies (PIC) improve
                      learning efficiency vs.\ representational alternatives? \\
    \textbf{RQ3.2} & How robust are graph-based policies to sensor and
                      composition changes? \\
    \textbf{RQ3.3} & What are the computational costs of architectural
                      vs.\ representational invariance? \\
    \bottomrule
\end{tabular}
\end{center}
    \note{Walk through at a high level — the committee has read the prospectus.
    A good framing: ``Eight research questions, three clusters, one thread.''
    Gesture down the table by group (about 20 seconds per group) and move on.
    Total time target: 1 minute.}
\end{frame}


\begin{frame}[shrink=24]{Research Objectives}
\vspace{1em}
\begin{center}
\renewcommand{\arraystretch}{1.2}
\small
\begin{tabular}{@{} l >{\raggedright\arraybackslash}p{0.75\linewidth} @{}}
    \toprule
    \multicolumn{2}{@{}l}{\textbf{C1 — Policy Upsampling}} \\
    \midrule
    \textbf{RO1.1} & Design curriculum-based team-scaling via policy duplication and retraining \\
    \textbf{RO1.2} & Define and validate agent-steps as a normalized training cost metric \\
    \textbf{RO1.3} & Evaluate curriculum effectiveness across distinct heterogeneity environments \\
    \textbf{RO1.4} & Identify task characteristics that moderate curriculum benefit or failure \\
    \midrule
    \multicolumn{2}{@{}l}{\textbf{C2 — Implicit Indication}} \\
    \midrule
    \textbf{RO2.1} & Formalize implicit indication as a representational approach to parameter sharing \\
    \textbf{RO2.2} & Design HyperGrid environment for controlled heterogeneity experiments \\
    \textbf{RO2.3} & Evaluate learning efficiency vs.\ per-type baselines \\
    \textbf{RO2.4} & Characterize robustness and zero-shot generalization \\
    \textbf{RO2.5} & Assess computational and storage costs \\
    \midrule
    \multicolumn{2}{@{}l}{\textbf{C3 — Paradigm Comparison}} \\
    \midrule
    \textbf{RO3.1} & Implement PIC and train under matched experimental conditions \\
    \textbf{RO3.2} & Compare learning efficiency: PIC vs.\ implicit indication \\
    \textbf{RO3.3} & Evaluate relative robustness to sensor and composition changes \\
    \textbf{RO3.4} & Assess computational costs of each paradigm \\
    \textbf{RO3.5} & Identify domain characteristics favoring each approach \\
    \bottomrule
\end{tabular}
\end{center}

    \note{Walk through by contribution group — about 15 seconds per group.
    The committee is checking coverage, not evaluating individual RO wording.
    A simple framing: ``Fourteen objectives across three contributions.
    Everything on this slide was executed; the next slide shows the mapping.''
    Total time target: 1 minute.}
\end{frame}


\begin{frame}[shrink=20]{Research Questions and Objectives}
\vspace{1.5em}
\begin{center}
\renewcommand{\arraystretch}{1.35}
\small
\begin{tabular}{@{} l c l @{}}
    \toprule
    \textbf{Question} & & \textbf{Objectives} \\
    \midrule
    \multicolumn{3}{@{}l}{\textbf{C1 — Policy Upsampling} \hfill \textit{Ch.\ 3}} \\
    \midrule
    RQ1.1 & $\longrightarrow$ & RO1.1, RO1.2, RO1.3 \\
    RQ1.2 & $\longrightarrow$ & RO1.3, RO1.4 \\
    \midrule
    \multicolumn{3}{@{}l}{\textbf{C2 — Implicit Indication} \hfill \textit{Ch.\ 4}} \\
    \midrule
    RQ2.1 & $\longrightarrow$ & RO2.1, RO2.2, RO2.3 \\
    RQ2.2 & $\longrightarrow$ & RO2.4 \\
    RQ2.3 & $\longrightarrow$ & RO2.5 \\
    \midrule
    \multicolumn{3}{@{}l}{\textbf{C3 — Paradigm Comparison} \hfill \textit{Ch.\ 5}} \\
    \midrule
    RQ3.1 & $\longrightarrow$ & RO3.1, RO3.2 \\
    RQ3.2 & $\longrightarrow$ & RO3.3 \\
    RQ3.3 & $\longrightarrow$ & RO3.4, RO3.5 \\
    \bottomrule
\end{tabular}
\end{center}

    \note{This slide is a navigational reference — the committee is checking coverage.
    You do not need to read each row.
    A single sentence works: ``Every RQ maps to one or more objectives,
    all executed in the corresponding chapter.''
    Call out the shared objective if it catches their eye:
    ``RO1.3 appears twice because evaluating across environments
    answers both the efficiency and the heterogeneity question.''
    Total time target: 30--45 seconds.}
\end{frame}

